\section{Earlier comparisons and full final metrics}\label{app:earlier}

\subsection{E1: original recovery comparison}
E1 comprises 216 episodes and 13 methods, with 237,600 eligible points and 135 intervention events. Twenty-seven events have missing command context. Nineteen clean episodes provide 20,900 eligible calibration points; the principal state-machine threshold search allows at most 209 alarmed points on this finite population. This differs from the final audit's point-quantile calibration. The additional constant episode band uses rank 19 of 19 episode maxima; shifted tasks do not establish exchangeability or conformal coverage. Matched-entry ablations preserve their own realized alarm burdens.
\begin{table*}[t]\centering
\caption{E1 attack and recovery outcomes, retained from the saved presentation. Clearance counts condition on detected events with measured return.}
\setlength{\tabcolsep}{3pt}\begin{tabular}{lrrrrr}\toprule
Method & Benign & F1 & Overlap/fresh & Early & Clear/return \\\midrule
Persistence & 878 & 0.057 & 127/127 & 127 & 127/127 \\
ARX one-step & 792 & 0.803 & 108/108 & 108 & 108/108 \\
Following point & 794 & 0.808 & 108/108 & 108 & 108/108 \\
Following hysteresis & 993 & 0.785 & 108/108 & 108 & 108/108 \\
Nominal hysteresis & 789 & 0.713 & 108/108 & 45 & 108/108 \\
Nominal fixed return & 621 & 0.657 & 108/108 & 0 & 108/108 \\
Learned point & 819 & 0.700 & 108/108 & 18 & 108/108 \\
Learned hysteresis & 792 & 0.683 & 108/108 & 0 & 108/108 \\
Learned fixed return & 827 & 0.690 & 108/108 & 0 & 108/108 \\
Uncertainty return & 309 & 0.446 & 108/104 & 0 & 36/108 \\
Episode CP band & 0 & 0.744 & 108/108 & 75 & 108/108 \\
Matched no allowance & 140 & 0.709 & 108/108 & 0 & 108/108 \\
Matched no dwell & 20 & 0.453 & 108/108 & 0 & 36/108 \\
\bottomrule\end{tabular}

\end{table*}
\begin{table*}[t]\centering
\caption{E1 additional metrics. Unknown states remain eligible, cannot count as detections, and receive AP score $-1$, below valid nonnegative scores.}
\setlength{\tabcolsep}{3pt}\begin{tabular}{lrrrrrr}\toprule
Method & Precision & Recall & AP & Episodes/h & Missed & Known (\%) \\\midrule
Persistence & 0.203 & 0.033 & 0.130 & 480.8 & 8 & 100.0 \\
ARX one-step & 0.868 & 0.748 & 0.770 & 1426.3 & 27 & 90.0 \\
Following point & 0.878 & 0.747 & 0.770 & 793.9 & 27 & 90.0 \\
Following hysteresis & 0.841 & 0.735 & 0.770 & 99.0 & 27 & 90.0 \\
Nominal hysteresis & 0.706 & 0.719 & 0.587 & 84.8 & 27 & 90.0 \\
Nominal fixed return & 0.656 & 0.658 & 0.587 & 48.5 & 27 & 90.0 \\
Learned point & 0.658 & 0.748 & 0.589 & 1054.5 & 27 & 90.0 \\
Learned hysteresis & 0.628 & 0.747 & 0.589 & 62.6 & 27 & 90.0 \\
Learned fixed return & 0.641 & 0.747 & 0.589 & 115.2 & 27 & 90.0 \\
Uncertainty return & 0.318 & 0.747 & 0.589 & 26.3 & 27 & 90.0 \\
Episode CP band & 0.742 & 0.746 & 0.589 & 0.0 & 27 & 90.0 \\
Matched no allowance & 0.675 & 0.747 & 0.589 & 28.3 & 27 & 90.0 \\
Matched no dwell & 0.325 & 0.747 & 0.589 & 28.3 & 27 & 90.0 \\
\bottomrule\end{tabular}

\end{table*}
\subsection{E2: normal commissioning}
The finite grid contains ARX orders 4, 8, 16 and 32 and finite impulse response lengths 16, 32 and 64, each with ridge $10^{-7}$, $10^{-5}$ or $10^{-3}$, on three plants: 63 fits. Fitting and selection precede 81 new normal episodes. The selected orders are 8, 16 and 32 with ridge $10^{-7}$. Table~\ref{tab:normalall} retains distinct prediction roles. Forced-latch probes measure normal clearance feasibility, not post-attack recovery. The hard-plant reductions in empirical allowance do not establish a probabilistic or deterministic error bound.
\begin{table*}[t]\centering
\caption{E2 complete normal-reference comparison. Each row has 29,700 eligible observations and 108 forced-latch probes. Selected reference and selected ARX are identical and appear once.}\label{tab:normalall}
\setlength{\tabcolsep}{4pt}\begin{tabular}{llrrrr}\toprule
Plant & Reference & RMSE & Allowance & Feasible (\%) & Clear/probes \\\midrule
Q1 & ARX(4), free & 0.00948 & 0.0553 & 99.80 & 108/108 \\
Q1 & Selected ARX, free & 0.00934 & 0.0553 & 99.82 & 108/108 \\
Q1 & Selected FIR & 0.01211 & 0.2322 & 0.00 & 0/108 \\
Q1 & ARX(4), one-step & 0.00976 & 0.0649 & 97.92 & 108/108 \\
Q1 & Command as state & 0.09584 & 2.1344 & 0.00 & 0/108 \\
Q2 & ARX(4), free & 0.01020 & 0.1469 & 0.00 & 0/108 \\
Q2 & Selected ARX, free & 0.00972 & 0.0716 & 93.32 & 108/108 \\
Q2 & Selected FIR & 0.01693 & 0.3971 & 0.00 & 0/108 \\
Q2 & ARX(4), one-step & 0.01036 & 0.1211 & 0.00 & 0/108 \\
Q2 & Command as state & 0.10423 & 2.1528 & 0.00 & 0/108 \\
Q3 & ARX(4), free & 0.01586 & 0.1913 & 0.00 & 0/108 \\
Q3 & Selected ARX, free & 0.01439 & 0.1175 & 0.00 & 0/108 \\
Q3 & Selected FIR & 0.02215 & 0.4596 & 0.00 & 0/108 \\
Q3 & ARX(4), one-step & 0.01573 & 0.1807 & 0.00 & 0/108 \\
Q3 & Command as state & 0.11251 & 2.2058 & 0.00 & 0/108 \\
\bottomrule\end{tabular}

\end{table*}
\subsection{E5: complete additional metrics and plant strata}
The following tables complement the complete 17-method audit in Table~\ref{tab:audit}. Alarm starts per hour use the total eligible duration of wholly benign episodes. Their high frequency for point rules reflects repeated short excursions; the metric should be read alongside occupancy. No unknown-state denominator trimming is applied to this audit. All 189 episodes and 3,213 outcomes are retained.
\begin{table*}[t]\centering
\caption{Final audit: unadjusted point precision/recall, benign occupancy, alarm starts per benign hour and missed events.}
\begin{tabular}{lrrrrr}\toprule
Method & Precision & Recall & Benign (\%) & Starts/h & Misses \\\midrule
base point & 0.6434 & 0.9931 & 0.957 & 1537.4 & 0 \\
base return & 0.4777 & 0.9936 & 7.644 & 1208.1 & 0 \\
base long & 0.4699 & 0.9936 & 8.031 & 1208.1 & 0 \\
base dual & 0.4699 & 0.9936 & 8.031 & 1208.1 & 0 \\
linear point & 0.6421 & 0.9930 & 1.004 & 1618.2 & 0 \\
linear dual & 0.4651 & 0.9934 & 8.415 & 1270.7 & 0 \\
RBF point & 0.6445 & 0.9929 & 0.937 & 1501.0 & 0 \\
RBF return & 0.4811 & 0.9932 & 7.466 & 1191.9 & 0 \\
RBF long & 0.4729 & 0.9932 & 7.870 & 1191.9 & 0 \\
RBF dual & 0.4729 & 0.9932 & 7.870 & 1191.9 & 0 \\
RBF dual, no dwell & 0.6214 & 0.9929 & 1.341 & 1501.0 & 0 \\
following point & 0.8575 & 0.9930 & 0.973 & 1618.2 & 0 \\
following dual & 0.5487 & 0.9934 & 8.569 & 1331.3 & 0 \\
RBF following dual & 0.5486 & 0.9934 & 8.617 & 1317.2 & 0 \\
persistence & 0.1286 & 0.0192 & 0.591 & 274.7 & 41 \\
mean (21) & 0.0385 & 0.0046 & 0.094 & 78.8 & 60 \\
base CUSUM & 0.2665 & 0.9887 & 0.962 & 301.0 & 0 \\
\bottomrule\end{tabular}

\end{table*}
\begin{table*}[t]\centering
\caption{Primary base-dual versus RBF-dual comparison by plant. Every plant has 29,700 benign points and 36 attack events.}
\setlength{\tabcolsep}{4pt}\begin{tabular}{lrrrrrr}\toprule
Plant & Base benign & RBF benign & Base gaps & RBF gaps & Base fresh & RBF fresh \\\midrule
Q1 & 2397 & 2522 & 0 & 0 & 36 & 36 \\
Q2 & 2646 & 2532 & 2 & 3 & 29 & 33 \\
Q3 & 2113 & 1958 & 0 & 0 & 36 & 36 \\
\bottomrule\end{tabular}

\end{table*}
